\documentclass[11pt]{article}

\usepackage[margin=0.9in]{geometry}
\usepackage{amsmath,amssymb}
\usepackage{booktabs}
\usepackage{enumitem}
\usepackage{listings}
\usepackage{xcolor}
\usepackage{hyperref}

\definecolor{backcolour}{rgb}{0.95,0.95,0.92}
\definecolor{codeblue}{rgb}{0.05,0.15,0.55}
\definecolor{codegreen}{rgb}{0.0,0.45,0.0}

\lstset{
    backgroundcolor=\color{backcolour},
    basicstyle=\ttfamily\small,
    keywordstyle=\color{codeblue}\bfseries,
    commentstyle=\color{codegreen},
    stringstyle=\color{purple},
    breaklines=true,
    frame=single,
    showstringspaces=false,
    tabsize=4
}

\title{CPSC 250 \\ Class Customization}
\author{Edward Brash}
\date{}

\begin{document}

\maketitle

\section{Introduction}

One of the most common things we want to do with an object is print information about it.

For example, suppose we have:

\begin{lstlisting}[language=Python]
moes = Restaurant("Moe's", "Mexican")
\end{lstlisting}

It would be convenient if we could simply write:

\begin{lstlisting}[language=Python]
print(moes)
\end{lstlisting}

and automatically obtain useful information about the object.

Python provides a special mechanism for doing exactly this.

\section{The \texttt{\_\_str\_\_()} Method}

Python classes can define a special method:

\begin{lstlisting}[language=Python]
def __str__(self):
\end{lstlisting}

This method tells Python how the object should appear when printed.

The method should return a string.

Example:

\begin{lstlisting}[language=Python]
def __str__(self):
    return f"{self.name} ({self.cuisine})"
\end{lstlisting}

\section{Automatic Behavior}

Suppose we define:

\begin{lstlisting}[language=Python]
class Restaurant:

    def __init__(self, name, cuisine):
        self.name = name
        self.cuisine = cuisine

    def __str__(self):
        return f"{self.name} ({self.cuisine})"
\end{lstlisting}

Then:

\begin{lstlisting}[language=Python]
moes = Restaurant("Moe's", "Mexican")

print(moes)
\end{lstlisting}

automatically calls:

\begin{lstlisting}[language=Python]
moes.__str__()
\end{lstlisting}

behind the scenes.

\section{Important Note}

The \verb|__str__| method must return a single string.

It does not directly print the information itself.

Bad design:

\begin{lstlisting}[language=Python]
def __str__(self):
    print(self.name)
\end{lstlisting}

Correct design:

\begin{lstlisting}[language=Python]
def __str__(self):
    return self.name
\end{lstlisting}

The returned string is then sent to the \verb|print()| statement.

\section{Benefits of \texttt{\_\_str\_\_()}}

The \verb|__str__| method makes objects easier to work with.

Advantages include:

\begin{itemize}
    \item cleaner debugging,
    \item easier printing,
    \item more readable output,
    \item more professional object design.
\end{itemize}

\section{Example}

Suppose we define:

\begin{lstlisting}[language=Python]
class Restaurant:

    def __init__(self, name, cuisine):
        self.name = name
        self.cuisine = cuisine

    def __str__(self):
        return f"Restaurant: {self.name}, Cuisine: {self.cuisine}"
\end{lstlisting}

Then:

\begin{lstlisting}[language=Python]
moes = Restaurant("Moe's", "Mexican")

print(moes)
\end{lstlisting}

produces:

\begin{lstlisting}
Restaurant: Moe's, Cuisine: Mexican
\end{lstlisting}

\section{Customizing Objects}

This is our first example of \textbf{class customization}.

Python allows us to redefine how many built-in operations behave for our own classes.

This leads naturally to operator overloading.

\section{Transition to Operator Overloading}

A natural question is:

\begin{quote}
What does it mean to add two objects together?
\end{quote}

For example:

\begin{lstlisting}[language=Python]
moes = Restaurant()
panera = Restaurant()

fusion = moes + panera
\end{lstlisting}

What should that mean?

The answer depends on the class.

For our own classes, we can define the rules ourselves.

\section{Algebra and Operators}

Algebra is really a set of rules for manipulating symbols and objects.

Operators include:

\begin{itemize}
    \item arithmetic operators:
    \[
    +, -, *, /, \%
    \]

    \item comparison operators:
    \[
    <, >, <=, >=, ==, !=
    \]
\end{itemize}

Python allows us to redefine these operators for our own classes.

\section{Example: A Time Class}

Suppose we define a \verb|Time| class:

\begin{lstlisting}[language=Python]
class Time:

    def __init__(self, hours, minutes):
        self.hours = hours
        self.minutes = minutes

    def __str__(self):
        return f"{self.hours}:{self.minutes}"
\end{lstlisting}

Now we can print time objects naturally.

\begin{lstlisting}[language=Python]
t1 = Time(10, 40)

print(t1)
\end{lstlisting}

\section{Comparison Operators}

Suppose we want to compare two time objects.

For example:

\begin{lstlisting}[language=Python]
time1 < time2
\end{lstlisting}

Python does not automatically know how to compare two custom objects.

We must define the comparison ourselves.

\section{The \texttt{\_\_lt\_\_()} Method}

The special method:

\begin{lstlisting}[language=Python]
def __lt__(self, other):
\end{lstlisting}

defines the meaning of the less-than operator:

\[
<
\]

for objects of the class.

\section{Example Implementation}

\begin{lstlisting}[language=Python]
class Time:

    def __init__(self, hours, minutes):
        self.hours = hours
        self.minutes = minutes

    def __str__(self):
        return f"{self.hours}:{self.minutes}"

    def __lt__(self, other):

        if self.hours < other.hours:
            return True

        elif self.hours == other.hours:

            if self.minutes < other.minutes:
                return True

        return False
\end{lstlisting}

\section{Using the Comparison}

Now suppose:

\begin{lstlisting}[language=Python]
time1 = Time(10, 40)
time2 = Time(12, 15)
time3 = Time(9, 15)
\end{lstlisting}

We can search for the earliest time:

\begin{lstlisting}[language=Python]
min_time = time1

if time2 < min_time:
    min_time = time2

if time3 < min_time:
    min_time = time3

print(f"Earliest time is {min_time}")
\end{lstlisting}

Because we defined \verb|__lt__|, Python now knows how to compare Time objects.

\section{Rich Comparison Operators}

Python provides several comparison customization methods.

\begin{center}
\begin{tabular}{cc}
\toprule
Method & Operator \\
\midrule
\verb|__lt__| & < \\
\verb|__le__| & <= \\
\verb|__gt__| & > \\
\verb|__ge__| & >= \\
\verb|__eq__| & == \\
\verb|__ne__| & != \\
\bottomrule
\end{tabular}
\end{center}

\section{Important Warning}

Some comparison methods have default behavior, but that behavior may not match what we actually want.

Therefore:

\begin{quote}
Be careful and define comparison behavior explicitly whenever meaningful comparisons are needed.
\end{quote}

\section{Key Takeaways}

\begin{itemize}
    \item The \verb|__str__| method controls how objects print.
    \item \verb|print(object)| automatically calls \verb|object.__str__()|.
    \item Operator overloading allows us to redefine operators for custom classes.
    \item Comparison operators are implemented using special methods such as \verb|__lt__|.
    \item Different classes may define operators in different ways.
\end{itemize}

\end{document}
